\documentclass[a4paper,slidestop,xcolor=pst,blue]{beamer}

\usepackage{beamerthemesplit}
\usepackage[utf8]{inputenc}
\usepackage[spanish]{babel}
\usepackage{graphicx}
\usepackage{pstricks} % PSTricks package
\usepackage{setspace}
\usepackage{multirow}
\usepackage{listings}
\usepackage{pgfpages}
\usepackage{hyperref}
\usepackage{etoolbox}
\usepackage{epstopdf}

\makeatletter
\patchcmd{\beamer@sectionintoc}{\vskip1.5em}{\vskip0.5em}{}{}
\makeatother

\setbeamercovered{dynamic}
\setcounter{tocdepth}{2}
\setbeamercolor{frametitle}{fg=black,bg=white}
\setbeamercolor{section in toc shaded}{fg=black}
\setbeamercolor{section in toc}{fg=red}
\setbeamercolor{subsection in toc shaded}{fg=black}
\setbeamercolor{subsection in toc}{fg=red}
\setbeamerfont{section in toc}{size=\small}
\setbeamerfont{subsection in toc}{size=\small}
\setbeamertemplate{section in toc shaded}[default][99]
\setbeamertemplate{subsection in toc shaded}[default][99]

\AtBeginSection[]
{\begin{frame}[c]
  \frametitle{Índice}
	\tableofcontents[currentsection,
        sectionstyle=show/shaded,
        subsectionstyle=hide]
\end{frame}}

\AtBeginSubsection[]
{\begin{frame}[c]
	\frametitle{Índice}
	\tableofcontents[
  		currentsection,
  		sectionstyle=shaded/shaded,
  		currentsubsection,
  		subsectionstyle=show/shaded/hide
		]
\end{frame}}

\setbeamercolor{frametitle}{fg=black,bg=white}

\setbeamertemplate{frametitle}{
	\begin{centering}
		\insertframetitle
		\par
	\end{centering}
}

\usetheme[secheader]{Boadilla} 

\title[Procesos de Negocio]{Especificación y Modelado de Procesos de Negocio}

\author[P. S{\'a}nchez]{\alert{Pablo S{\'a}nchez}}

\institute[IIE]{
		   Dpto. Ingenier{\'i}a Inform{\'a}tica y Electr{\'o}nica \\
		   Universidad de Cantabria \\
		   Santander (Cantabria, Espa{\~n}a) \\
		   \texttt{p.sanchez@unican.es}
}

\date{}

\begin{document}

\begin{frame}[c]
	\titlepage
	\begin{columns}
		\column{0.50\linewidth}
			\centering
    		\includegraphics[width=.28\textwidth,keepaspectratio=true]{images/istr.eps}
		\column{0.50\linewidth}
			\centering
			\includegraphics[width=.25\textwidth,keepaspectratio=true]{images/uc.eps}
	\end{columns}
\end{frame}

\begin{frame}[c]
    \frametitle{\alert{Advertencia}}
    \begin{center}
        Todo el material contenido en este documento no constituye en modo alguno una obra de referencia o apuntes oficiales mediante el cual se puedan preparar las pruebas evaluables necesarias para superar la asignatura. \ \\
        \ \\
        Este documento contiene exclusivamente una serie de diapositivas cuyo objetivo es servir de complemento visual a las actividades realizadas en el aula para la transmisi{\'o}n del contenido sobre el cual versar{\'a}n las mencionadas pruebas evaluables.  \ \\
        \ \\
        Dicho de forma m{\'a}s clara, \alert{estas transparencias no son apuntes y su objetivo no es servir para que el alumno pueda preparar la asignatura.}
    \end{center}
\end{frame}

\section{Introducción}

\subsection{Objetivos y Bibliografía}

\begin{frame}[c]
    \frametitle{Objetivos del Tema}
    \begin{enumerate}[<+->]
         \item Comprender qué es un proceso de negocio.
         \item Conocer qué es BPMN 2.0.
         \item Comprender cómo se utilizan los principales elementos de BPMN 2.0.
         \item Ser capaz de modelar procesos de negocio básicos en BPMN 2.0.
    \end{enumerate}
\end{frame}

\begin{frame}[c]
    \frametitle{Bibliografía}
    \begin{thebibliography}{1}
        \bibitem[\protect\astroncite{OMG}{2011}]{bpmn:2011}
        Object Management Group (OMG). (2011).
        \newblock {\em {Business Process Model and Notation (BPMN) 2.0}}.
        \newblock OMG Document Number: formal/2011-01-03.
    \end{thebibliography}
\end{frame}

\section{Proceso de Negocio}

\begin{frame}[c]
    \frametitle{Proceso de Negocio}
    \begin{block}{Proceso de Negocio~\cite{}}
        Un \emph{proceso de negocio} es una colección de actividades que aceptan una o más entradas y producen una salida que es de valor para un cliente. Cada proceso de negocio tiene un objetivo, el cual se ve influenciado por eventos procedentes del exterior o de otros procesos.
    \end{block}
    \uncover<2->{
        \begin{block}{Proceso de Negocio~\cite{}}
            Un \emph{proceso de negocio} es una conjunto estructurado y medible de actividades diseñados para producir una salida concreta para un cliente o mercado. Un proceso es por tanto un conjunto de actividades ordenados en el tiempo y el espacio, con un principio, un fin y entradas y salidas claramente identificadas.
        \end{block}
    }
\end{frame}

\begin{frame}[c]
    \frametitle{Características de un Proceso de Negocio}
    \begin{enumerate}[<+->]
        \item Contiene un conjunto finito de actividades.
        \item Sus actividades están parcialmente ordenadas.
        \item Cada actividad tiene asociado uno o más eventos que originan su ejecución.
        \item Cada actividad produce una salida.
        \item La salida general del proceso es de interés para un consumidor.
    \end{enumerate}
\end{frame}

\begin{frame}[c]
    \frametitle{Ejemplo de Proceso de Negocio}
    \begin{center}
        \href{http://youtu.be/gQ_OHOelPvM}{Ejemplo Informal I} \ \\
        \href{https://youtu.be/bfB38RlXRbI}{Ejemplo Informal II} \ \\
        \href{https://youtu.be/vvbN-cWe0A0}{Ejemplo Informal II} 
    \end{center}
\end{frame}

\section{BPMN 2.0}

\subsection{Introducción}

\begin{frame}[c]
    \frametitle{Business Process Modelling Notation (BPMN)}
    \begin{block}{BPMN 2.0}
        \emph{Business Process Modelling Notation} (BPMN) es un lenguaje de modelado de procesos de negocio estandarizado por la OMG (\emph{Object Managament Group}) con el objetivo de proporcionar una notación fácilmente entendible por todos los actores de un proceso de negocio, abarcando desde los analistas que lo definen hasta los programadores que desarrollan las herramientas que dan soporte al mismo.
    \end{block}
\end{frame}

\begin{frame}[c]
    \frametitle{Características BPMN}
    \begin{enumerate}[<+->]
        \item Intento de proporcionar una notación de alto nivel para WSBPEL.
        \item Notación basada en diagramas de flujo.
        \item Semántica basada en tokens.
    \end{enumerate}
\end{frame}

\subsection{Tipos de Procesos}

\begin{frame}[c]
    \frametitle{Tipos de Procesos de Negocio}
    \begin{description}[<+->]
        \item[Privados] Internos a una organización.
        \item[Públicos] Interacción entre un proceso privado y otro proceso o un participante.
        \item[Colaboración] Interacción entre dos o más entidades de negocio.
        \item[Coreografías] Contrato o definición del comportamiento de la interacción entre varios participantes.
        \item[Conversaciones] Intercambio de mensajes y objetos de negocio entre entidades.
    \end{description}
\end{frame}

%%\begin{frame}[c]
%%    \frametitle{Elementos de BPMN 2.0}
%%    \begin{enumerate}[<+->]
%%        \item Flow Objects.
%%        \item Data.
%%        \item Swimlanes.
%%        \item Artifacts.
%%    \end{enumerate}
%%\end{frame}

\subsection{Pools and Lanes}

\begin{frame}[c]
    \begin{block}{Pool}
        Un \emph{pool} representa un participante u organización dentro de un proceso de negocio.
        \ \\
        \begin{center}
            \includegraphics[width=0.25\textwidth]{symbols/pool.eps}
        \end{center}
    \end{block}
    \begin{block}{Lane}
        Una \emph{pool} representa una entidad perteneciente a una determinada organización y que realiza una serie de tareas concretas dentro de un proceso de negocio.  
        \ \\
        \begin{center}
            \includegraphics[width=0.25\textwidth]{symbols/lane.eps}
        \end{center}
    \end{block}
\end{frame}

\subsection{Eventos, Tareas y Flujos}

\begin{frame}[c]
    \frametitle{Eventos y Tareas}
    \begin{block}{Evento}
        Un \alert{\emph{evento}} es algo que ocurre durante el curso de un proceso de negocio. \uncover<2->{Un evento sucede a consecuencia de una causa (\emph{trigger}) y puede poseer un impacto (\emph{result}).}
        \uncover<3->{Los eventos puede ser de \emph{inicio}, \emph{intermedios} o de \emph{finalización}.}
        \begin{columns}
        \begin{column}{width=0.33\textwidth}
            \begin{center}
                \includegraphics[width=0.50\textwidth]{symbols/startEvent.eps}
            \end{center}
        \end{column}
        \begin{column}{width=0.33\textwidth}
            \begin{center}
                \includegraphics[width=0.50\textwidth]{symbols/intermediateEvent.eps}
            \end{center}
        \end{column}
        \begin{column}{width=0.33\textwidth}
            \begin{center}
                \includegraphics[width=0.50\textwidth]{symbols/endEvent.eps}
            \end{center}
        \end{column}
    \end{block}
    \begin{block}{Tarea}
        Una \alert{\emph{tarea}} es una unidad de trabajo atómica realizada por una entidad dentro de un proceso de negocio.
        \ \\
        \begin{center}
            \includegraphics[width=0.25\textwidth]{symbols/task.eps}
        \end{center}
    \end{block}
\end{frame}

\begin{frame}[c]
    \frametitle{Flujos de Control}
    \begin{block}{Sequence Flow}
        Un \emph{flujo de secuencia} especifica el orden en el cual se han de ejecutar las actividades de un proceso de negocio. Representa la transferencia de un token de una actividad de origen a una actividad de destino. 
        \ \\
        \begin{center}
            \includegraphics[width=0.5/textwidth]{symbols/sequenceFlow.eps}
        \end{center}
    \end{block}
    \begin{block}{Message Flow}
        Un \emph{flujo de mensajes} especifica cómo dos participantes de un proceso de negocio intercambian un mensaje. 
        \ \\
        \begin{center}
            \includegraphics[width=0.5/textwidth]{symbols/messageFlow.eps}
        \end{center}
    \end{block}
\end{frame}

\begin{frame}[c]
    \frametitle{Tipos de Eventos Principales}
    \begin{center}
        \includegraphics[width=\linewidth]{symbols/eventKinds.eps}
    \end{center}
\end{frame}

\subsection{Ramificaciones y Uniones}

\begin{frame}[c]
    \frametitle{Ramificaciones y Uniones}
    \begin{block}{Gateways}
        Un \emph{gateway} es una estructura que se utiliza para dividir o unir flujos de ejecución.
        \ \\
        \begin{center}
            \includegraphics[width=0.25\text]{symbols/gateway.eps}
        \end{center}
    \end{block}
    \begin{block}{Conditional Flow}
        Un \emph{flujo condicional} es un flujo con una condición asociada que se evalúa en tiempo de ejecución para determinar cuando el flujo estará activo y por tanto fluirá un token a través de él.
        \ \\
        \begin{center}
            \includegraphics[width=0.5/textwidth]{symbols/conditionalFlow.eps}
        \end{center}
    \end{block}
    \begin{block}{Default Flow}
        Un \emph{flujo por defecto} es un flujo condicional cuya guarda es siempre verdadera.
        \ \\
        \begin{center}
            \includegraphics[width=0.5/textwidth]{symbols/defaultFlow.eps}
        \end{center}
    \end{block}
\end{frame}

\begin{frame}[c]
    \frametitle{Tipos de Gateway}
    \begin{center}
            \includegraphics[width=\linewidth]{symbols/gatewayKinds.eps}
    \end{center}
\end{frame}

\section{Sumario}

\begin{frame}[c]
    \frametitle{¿Qué tengo que saber de todo esto?}
    \begin{enumerate}[<+->]
        \item Entender el concepto de proceso de negocio.
        \item Conocer los diferentes tipos de proceso de negocio que existen.
        \item Conocer la utilidad y filosofía de BPMN 2.0.
        \item Ser capaz de modelar procesos de negocio sencillos utilizando BPMN 2.0.
    \end{enumerate}
\end{frame}

\end{document}
